\documentclass[10pt,twocolumn]{article}

% --- Page and font ---
\usepackage[letterpaper,margin=0.50in]{geometry}
\usepackage{mathpazo}
\usepackage[T1]{fontenc}
\usepackage{microtype}
\usepackage{amsmath,amssymb}
\usepackage{xcolor}
\usepackage{booktabs}
\usepackage{tabularx}
\usepackage{array}
\usepackage{ragged2e}
\usepackage{titlesec}
\usepackage{caption}
\usepackage{enumitem}
\usepackage{tcolorbox}
\usepackage[colorlinks=true,linkcolor=deepblue,citecolor=deepblue,urlcolor=deepblue]{hyperref}

% --- Muted palette matching the CVT appendix ---
\definecolor{deepblue}{RGB}{59,76,95}
\definecolor{lineblue}{RGB}{83,103,122}
\definecolor{softblue}{RGB}{235,242,247}
\definecolor{cellblue}{RGB}{221,232,238}
\definecolor{cellgreen}{RGB}{226,237,229}
\definecolor{cellgray}{RGB}{238,239,239}
\definecolor{mutedred}{RGB}{155,54,62}
\definecolor{mutedteal}{RGB}{55,115,120}

% --- Compact readable layout ---
\setlength{\columnsep}{0.28in}
\setlength{\parindent}{0pt}
\setlength{\parskip}{1.0pt}
\setlength{\textfloatsep}{3pt plus 1pt minus 1pt}
\setlength{\floatsep}{3pt plus 1pt minus 1pt}
\setlength{\intextsep}{3pt plus 1pt minus 1pt}
\captionsetup{font=scriptsize,labelfont=bf,skip=1pt}
\titleformat{\section}{\normalsize\bfseries\color{deepblue}}{\thesection}{0.45em}{}
\titlespacing*{\section}{0pt}{3.5pt}{1pt}
\renewcommand{\thesection}{\Alph{section}}
\pagestyle{plain}

\tcbset{
  cvtbox/.style={
    colback=softblue,
    colframe=lineblue,
    boxrule=0.45pt,
    arc=2pt,
    left=3pt,right=3pt,top=3pt,bottom=3pt,
    before skip=2pt,after skip=2pt
  },
  tablebox/.style={
    colback=white,
    colframe=lineblue,
    boxrule=0.35pt,
    arc=2pt,
    left=2pt,right=2pt,top=2pt,bottom=2pt,
    before skip=2pt,after skip=2pt
  }
}

\newcolumntype{Y}{>{\RaggedRight\arraybackslash}X}
\newcolumntype{W}[1]{>{\RaggedRight\arraybackslash}p{#1}}
\newcolumntype{M}[1]{>{\Centering\arraybackslash}p{#1}}
\newcommand{\code}[1]{\texttt{\fontsize{8.1}{9.3}\selectfont #1}}

\newenvironment{hptable}{%
  \begin{tcolorbox}[tablebox]
  \footnotesize
  \renewcommand{\arraystretch}{1.12}
  \setlength{\tabcolsep}{3.0pt}
}{%
  \end{tcolorbox}
}

\begin{document}

\appendix
\twocolumn[{
\begin{center}
{\LARGE\bfseries\color{deepblue} Appendix E: BERT Search Hyperparameters}\\[-0.2ex]
\end{center}
}]

\section{Core Algorithm}
These are the default high-level settings used by the encoder structured-pruning search.

\begin{hptable}
\begin{tabularx}{\linewidth}{@{}W{0.42\linewidth}Y@{}}
\toprule
\textbf{Parameter} & \textbf{Value} \\
\midrule
Population size & 120 \\
Number of generations & 80 \\
Elite size & 15 \\
Number of islands & 1 \\
\bottomrule
\end{tabularx}
\end{hptable}

\section{Archives}
The archives preserve high-quality candidates and prevent repeated macro signatures from dominating the stored finalists.

\begin{hptable}
\begin{tabularx}{\linewidth}{@{}W{0.52\linewidth}Y@{}}
\toprule
\textbf{Parameter} & \textbf{Value} \\
\midrule
Hall of Fame capacity & 50 \\
Hall of Fame max per macro signature & 3 \\
Pareto archive size & 100 \\
\bottomrule
\end{tabularx}
\end{hptable}

\section{Elites and Offspring}
Elitism preserves the strongest candidates, while the offspring schedule shifts toward larger edits as stagnation increases.

\begin{hptable}
\begin{tabularx}{\linewidth}{@{}W{0.43\linewidth}Y@{}}
\toprule
\textbf{Parameter} & \textbf{Value} \\
\midrule
Elite size & 15 \\
Per-island elites & 15 \\
Micro/macro split & 0.35 initial; 0.80 maximum; cap reached at stagnation 15 \\
Crossover rate & 0.25 micro, 0.65 macro, 0.40 default \\
\bottomrule
\end{tabularx}
\end{hptable}

\section{Stagnation}
Stagnation increases the probability of broader structural edits and can trigger additional full evaluations.

\begin{hptable}
\begin{tabularx}{\linewidth}{@{}W{0.45\linewidth}Y@{}}
\toprule
\textbf{Parameter} & \textbf{Value} \\
\midrule
Restart threshold & stagnation $\geq 20$ \\
Extra full evals & 24 \\
\bottomrule
\end{tabularx}
\end{hptable}

\section{EXP3-IX Bandit}
Each operator starts from a positive weight. The sampling probability mixes the learned distribution with uniform exploration.

\begin{hptable}
\begin{tabularx}{\linewidth}{@{}W{0.37\linewidth}Y@{}}
\toprule
\textbf{Parameter} & \textbf{Value} \\
\midrule
Exploration coefficient $\gamma$ & 0.07 \\
Initial weights & $w[o]=1.0$ for all operators \\
Probability & $(1-\gamma)\,w[o]/\sum_j w[j]+\gamma/K$ \\
Update & $w[o]\leftarrow w[o]\exp\!\left(\gamma\,\frac{r}{p+\gamma}\frac{1}{K}\right)$ \\
Overflow clip & normalise when $\max_o w[o] > 10^6$ \\
\bottomrule
\end{tabularx}
\end{hptable}

\section{Island Operator Priors}
Each island corresponds to a broad allocation style. The table below each label lists only the operators whose initial weights differ from the default value $1.0$.

\textbf{Head-heavy, low-neuron.} This island favours moving capacity toward attention heads while keeping the feed-forward allocation small.

\begin{hptable}
\begin{tabularx}{\linewidth}{@{}Y M{0.18\linewidth}@{}}
\toprule
\textbf{Operator} & \textbf{Weight} \\
\midrule
\code{budget\_neutral\_swap\_neurons\_for\_head} & 3.0 \\
\code{budget\_neutral\_swap\_head} & 2.0 \\
\code{head\_neuron\_trade} & 1.5 \\
\code{layer\_rewire} & 1.0 \\
\code{jitter} & 1.0 \\
\code{equilibrium\_mutate} & 0.8 \\
\bottomrule
\end{tabularx}
\end{hptable}

\textbf{Head-heavy, mid-neuron.} This island still favours attention, but keeps enough neuron capacity for mixed head--neuron refinements.

\begin{hptable}
\begin{tabularx}{\linewidth}{@{}Y M{0.18\linewidth}@{}}
\toprule
\textbf{Operator} & \textbf{Weight} \\
\midrule
\code{budget\_neutral\_swap\_neurons\_for\_head} & 2.0 \\
\code{budget\_neutral\_swap\_head} & 1.8 \\
\code{budget\_neutral\_swap\_neurons} & 1.4 \\
\code{equilibrium\_mutate} & 1.4 \\
\code{head\_neuron\_trade} & 1.2 \\
\code{jitter} & 1.0 \\
\bottomrule
\end{tabularx}
\end{hptable}

\textbf{Balanced.} This island starts near an even attention and feed-forward allocation, so its priors favour local swaps and equilibrium-preserving moves.

\begin{hptable}
\begin{tabularx}{\linewidth}{@{}Y M{0.18\linewidth}@{}}
\toprule
\textbf{Operator} & \textbf{Weight} \\
\midrule
\code{budget\_neutral\_swap\_head} & 1.8 \\
\code{budget\_neutral\_swap\_neurons} & 1.8 \\
\code{equilibrium\_mutate} & 1.6 \\
\code{group\_reallocate} & 1.2 \\
\code{macro\_budget\_shift} & 1.2 \\
\code{jitter} & 1.0 \\
\bottomrule
\end{tabularx}
\end{hptable}

\textbf{Neuron-heavy.} This island favours keeping more feed-forward capacity and uses operators that move budget from heads into neuron blocks.

\begin{hptable}
\begin{tabularx}{\linewidth}{@{}Y M{0.18\linewidth}@{}}
\toprule
\textbf{Operator} & \textbf{Weight} \\
\midrule
\code{budget\_neutral\_swap\_head\_for\_neurons} & 3.0 \\
\code{budget\_neutral\_swap\_neurons} & 2.2 \\
\code{neuron\_rebalance} & 1.6 \\
\code{group\_reallocate} & 1.2 \\
\code{jitter} & 1.0 \\
\bottomrule
\end{tabularx}
\end{hptable}

\textbf{Layer-sparse.} This island explores architectures with fewer active layers and stronger depth-level edits.

\begin{hptable}
\begin{tabularx}{\linewidth}{@{}Y M{0.18\linewidth}@{}}
\toprule
\textbf{Operator} & \textbf{Weight} \\
\midrule
\code{drop\_layer} & 3.0 \\
\code{drop\_multi} & 2.5 \\
\code{layer\_rewire} & 2.0 \\
\code{macro\_budget\_shift} & 1.8 \\
\code{budget\_neutral\_swap\_head} & 1.2 \\
\code{budget\_neutral\_swap\_neurons\_for\_head} & 1.2 \\
\code{jitter} & 1.0 \\
\bottomrule
\end{tabularx}
\end{hptable}

\end{document}
